\section{Structure of a Report}
\label{sec:structure}

Most scientific writing in engineering follows a variant of the
\textbf{IMRaD} pattern: \emph{Introduction, Methods, Results, and
Discussion}. \Cref{tab:structure} shows how this maps onto a typical
CS/engineering report and what each part is \emph{for} --- not just what
it is called.

\renewcommand{\arraystretch}{1.4}

\begin{table}[htbp]
  \centering
  \caption{Typical section structure for a CS/engineering report and the
  question each section must answer.}
  \label{tab:structure}

  \begin{tabular}{@{}p{4cm}p{10.5cm}@{}}
    \toprule
    \textbf{Section} & \textbf{Question it must answer} \\
    \midrule

    Title & What is this, in one line? \\
    \hline
    Abstract & In 150--250 words: problem, approach, key result. \\
    \hline
    Introduction & Why does this problem matter, and what exactly did you do about it? \\
    \hline
    Background / Related Work & What does the reader need to know, and how does your work relate to existing solutions? \\
    \hline
    Design / Methodology & How does your system or experiment work, and why did you make these design choices? \\
    \hline
    Implementation & What did you actually build? (Only if relevant --- keep it high-level; code goes in an appendix or repository.) \\
    \hline
    Evaluation / Results & What did you measure, and what did you observe? (Facts only, no interpretation yet.) \\
    \hline
    Discussion & What do the results \emph{mean}? What are the limitations? \\
    \hline
    Conclusion & What is the one-paragraph takeaway, and what would you do next? \\
    \hline
    References & Who else's work did you build on or compare against? \\

    \bottomrule
  \end{tabular}
\end{table}

\subsection{Abstract}
The abstract is the most-read and least-understood part of a report. It
is \emph{not} an introduction and it is \emph{not} a table of contents in
prose form. A good abstract compresses the whole report into four moves:
(1) one sentence of context/problem, (2) one sentence of what you did,
(3) one or two sentences of what you found, (4) optionally, one sentence
on why it matters. Write the abstract \emph{last}, after the rest of the
report is stable.

\subsection{Introduction vs.\ Background}
Students often merge these, which weakens both. The \textbf{introduction}
motivates the problem and states your contribution in the first page ---
a reader who stops after the introduction should still know what you did
and why. The \textbf{background/related work} section is where you teach
the reader the prerequisites and position your work relative to existing
approaches. If a fact is needed to understand your method, it belongs in
background, not scattered later.

\begin{tipbox}
  A useful test: could you swap the introduction of two unrelated reports
  and have both still make sense? If yes, your introduction is too generic
  --- it needs to say what \emph{your} project specifically does.
\end{tipbox}

\subsection{Results vs.\ Discussion}
Keep these separate, even if only conceptually. \textbf{Results} report
what happened, ideally with a table or figure and minimal commentary.
\textbf{Discussion} is where you interpret: Was this expected? Does it
support your hypothesis? What are the threats to validity (measurement
noise, small sample, unrepresentative test data, hardware differences)?
Mixing the two makes it hard for a reader --- or a grader --- to tell
which sentences are evidence and which are your opinion.

\begin{exercisebox}
  Take the last report you wrote (for this course or another) and
  highlight every sentence in your results section that contains a
  judgement word (``surprisingly'', ``unfortunately'', ``clearly better'').
  Each one is arguably discussion, not results. How many did you find?
\end{exercisebox}
